\chapter{Preferences, housekeeping and troubleshooting}
\label{ch:prefs}

This chapter covers the settings that shape the program itself --- as opposed to
the job --- and then the practical business of fixing things when they go wrong:
where your data lives, what to keep, and what to do when something misbehaves.

\section{The Preferences window}

Open it with \key{Ctrl+,}, \menu{Edit>Settings>Preferences}, or the
\btn{Preferences} button. It is a notebook, and its tabs each collect settings
from a different part of the program:

\begin{center}\small
\begin{tabular}{@{}L{2.7cm}L{11.5cm}@{}}
\toprule
\textbf{Tab} & \textbf{What is in it} \\
\midrule
\tabname{General} & Units and language, plus the settings for \emph{Start},
\emph{Options}, \emph{Laser}, \emph{Optimisations} and \emph{Laser Settings}. At
the bottom, a saving panel with \btn{Save}, \btn{Export} and \btn{Import}. \\
\tabname{Input/Output} & How files are read and written: pixels per inch,
\emph{Images}, \emph{SVG-Features}, \emph{DXF} and the I/O page. \\
\tabname{Classification} & Colour-based assignment of elements to operations,
with presets \btn{Automatic} and \btn{Manual}. \\
\tabname{Operations} & Defaults applied when operations are created. \\
\tabname{GUI} & Behaviour of the interface itself --- the table below. \\
\tabname{Scene} & Element display, selection behaviour and snapping. \\
\tabname{Colors} & Every colour used on the scene, with \btn{Reset Colors to
Default} and \btn{Reset Colors to brighter defaults} (for dark themes), plus the
object-label colour. \\
\tabname{Ribbon} & The ribbon editor: pages, panels and buttons. \\
\tabname{Coordinate Space} & Visible only with \cmd{developer_mode} on; shows the
coordinate transforms. \\
\bottomrule
\end{tabular}
\end{center}

\noindent The settings you are most likely to want, with their defaults:

\begin{center}\small
\begin{tabular}{@{}L{5.6cm}L{1.8cm}L{6.8cm}@{}}
\toprule
\textbf{Label} & \textbf{Default} & \textbf{Effect} \\
\midrule
\btn{UI-Colours} & System & Force \emph{Dark} or \emph{Light} regardless of the
operating system's setting. \\
\btn{Single Instance} & on & Allow only one MeerK40t at a time. Turn it off when
you genuinely want two devices driven from two windows. \\
\btn{Soundfile} & (system) & The sound played by \cmd{beep} --- the end-of-job
alarm for machines you cannot see. \\
\btn{Autosave workspace} & on & Periodic save to \file{\_autosave.svg}. See the
warning below. \\
\emph{Frequency} & 5 minutes & How often the autosave runs (30\,s to
10\,min). \\
\btn{Optimize element display} & on & Do not draw invisible elements. \\
\btn{Simplify Effects} & on & Limit effect complexity on screen; \key{CapsLock}
temporarily disables it. \\
\btn{Remember main window position} & on & Reopen where you left it. \\
\btn{Save Window Positions} & on & Reopen windows where you left them. \\
\btn{Launch Spooler on Job Start} & on & Bring up the spooler when a job
starts. \\
\btn{MouseWheel Pan} & off & Wheel pans instead of zooming; \key{Ctrl}+wheel does
the other. \\
\btn{Invert MouseWheel Pan} / \btn{Invert MouseWheel Zoom} & off & Reverse the
directions if they feel wrong. \\
\emph{ToolTip delay} & 100\unit{ms} & How quickly tooltips appear. \\
\emph{Ribbon: hover before description} & 2000\unit{ms} & How long to hover a
ribbon button before its long description appears. \\
\btn{Long UI descriptions} & on & Show extended help in ribbons and context
menus. \\
\emph{Level} & 0 & Warning severity: 1 all warnings, 2 normal + critical, 3
critical only, 4 ignore all. \\
\btn{Enable Laser Arm} & off & Require the laser to be armed (\btn{Arm}) before
it will start a job. \\
\btn{Track changes and allow undo} & on & Undo system on or off (needs a
restart). \\
\emph{Levels of Undo-States} & 20 & How many undo steps to keep (3--250). \\
\bottomrule
\end{tabular}
\end{center}

\begin{warnbox}[Autosave is not a recovery system]
MeerK40t writes \file{\_autosave.svg} every few minutes, but it does not offer to
reopen it after a crash or a bad exit. If something goes wrong, the file is there
--- you have to open it yourself. Do not treat it as a substitute for
\key{Ctrl+S}.
\end{warnbox}

\section{Appearance, language and updates}

\begin{itemize}
  \item \textbf{Themes}: \btn{UI-Colours} switches between the system setting,
        dark and light. The scene has its own palette in \tabname{Colors}; if a
        dark theme makes the bed or the grid hard to see, the single-click fix is
        \btn{Reset Colors to brighter defaults}.
  \item \textbf{Language}: \menu{Languages} switches the interface language.
        Changing it requires a restart, and the dialog says so. The bundled set
        includes English, German, Spanish, French, Hungarian, Italian, Japanese,
        Dutch, Polish, both Portuguese variants, Russian, Turkish and Chinese.
  \item \textbf{Updates}: the updater's settings live in the \emph{Options} page:
        \emph{Action} (\emph{No, thank you}, \emph{Look for major releases},
        \emph{Look for major+beta releases}) and \emph{Frequency} (\emph{At every
        startup}, \emph{Once per day}, \emph{Once per week}). The same check can
        be run on demand from \menu{Help>Check for Updates}, or with
        \cmd{check\_for\_updates} in the console; the result appears in a window
        titled \emph{Update-Info}. If you use beta releases, expect beta
        problems.
  \item \textbf{About} (\menu{Help>About MeerK40t}) has four tabs: \emph{About},
        \emph{In Memory of David Olsen}, \emph{System-Information} and
        \emph{MeerK40t-Components}. The last two list the versions of Python,
        wxPython, numpy, Pillow, potrace, vtracer, ezdxf, pyusb, pyserial, OpenCV,
        clipper and numba that are actually loaded --- exactly what to copy into a
        bug report.
  \item \textbf{Tips} (\menu{Help>Tips \&\& Tricks}) shows a tip with optional
        image, \btn{Previous tip}, \btn{Next tip} and \btn{Try it out} (which
        loads an example, so be careful: it can change your design). \btn{Show
        tips at startup} and \btn{Automatically Update} control whether you see
        them and whether they are refreshed from the project wiki.
\end{itemize}

\section{Where your data actually lives}

Knowing this list turns several classes of problem into a two-minute fix.

\begin{center}\small
\begin{tabular}{@{}L{5.0cm}L{9.2cm}@{}}
\toprule
\textbf{File} & \textbf{Where and why} \\
\midrule
\file{MeerK40t.cfg} & The configuration, in the per-user configuration directory
(\file{\textasciitilde/.config/MeerK40t/} on Linux,
\file{\textasciitilde/Library/Application Support/MeerK40t/} on macOS,
\file{\%LOCALAPPDATA\%\textbackslash MeerK40t\textbackslash} on Windows).
Settings, devices, window layout, recent files. \\
\file{MeerK40t.bak}, \file{.ba1}\dots\file{.ba5} & The five previous versions of
that file, newest first. \\
\file{\_autosave.svg} & The periodic workspace save, in the program's working
directory. \\
\file{cmdhistory.log} & Your console command history (last fifty commands). \\
\file{MeerK40t-\emph{timestamp}.txt} & Crash logs: written to the current working
directory, or to the configuration directory if that is not writable. \\
\file{history.csv} & Job history exported from the spooler with \emph{Write
CSV}. \\
\file{tips.txt} & The local cache of tips shown by the Tips window. \\
\bottomrule
\end{tabular}
\end{center}

\noindent The console also writes useful things on demand:
\cmd{channel save console} records the console output to a timestamped file,
\cmd{setting\_export <dir>} copies your whole configuration out (and
\cmd{setting\_import <file>} reads it back), and the \btn{Export} button in the
preferences does the same thing from the GUI.

\section{Troubleshooting by symptom}

\subsection*{The window is broken, a pane is missing, or a dialog is off-screen}

\begin{enumerate}[label=\arabic*.,leftmargin=1.4em]
  \item \menu{Panes>Reset Panes} restores the default layout.
  \item \menu{Tools>Reset Windows} forgets stored window positions; the console
        version is \cmd{window reset *}.
  \item Still wrong? Close the program, move \file{MeerK40t.cfg} aside, and start
        again. You lose the layout, not your projects.
\end{enumerate}

\subsection*{Nothing happens when I press Start}

The program is usually refusing on purpose, and it says why:

\begin{itemize}
  \item \emph{``Nothing to burn. Operations need assigned shapes\dots''} --- there
        is no artwork attached to any enabled operation.
  \item \emph{``Some shapes are not assigned to any operation\dots Queue
        anyway?''} --- some artwork will be skipped. Fix it with
        \menu{Edit>Select all non-assigned} and the operation squares.
  \item \emph{``Connect your laser device first''} --- the device is not talking
        to the machine; Chapter~\ref{ch:connect} has the checklist.
  \item \btn{Enable Laser Arm} is on and the laser is not armed --- press
        \btn{Arm}, or turn the preference off.
  \item Nothing at all, not even a dialog: check the \panel{Console} pane. The
        refusal is printed there.
\end{itemize}

\subsection*{A job stops in the middle, or the head goes somewhere unexpected}

\begin{enumerate}[label=\arabic*.,leftmargin=1.4em]
  \item Look at the \panel{Console} for the last device message; the driver's
        \btn{Controller} window carries the raw traffic and any alarm.
  \item An alarm (GRBL especially) usually means the machine hit a soft limit:
        check the bed size, the flips and the declared home corner
        (Chapter~\ref{ch:connect}) before assuming a mechanical fault.
  \item If the job is still queued, the \panel{Spooler} shows its status; a job
        can be paused, resumed, disabled or removed there.
  \item If the machine stopped but the program thinks it is running, use
        \btn{Abort} in the spooler and re-home before continuing.
\end{enumerate}

\subsection*{The burn is wrong}

Almost always one of four things, in this order of likelihood: focus is off; the
material is not what you think it is; the speed/power pair has not been verified
on this machine; or the raster settings (DPI, overscan, direction) are wrong for
the machine. Chapter~\ref{ch:tuning} is the systematic answer, and
Chapter~\ref{ch:engrave} covers raster-specific faults.

\subsection*{The program is slow}

\begin{itemize}
  \item Huge imported files: \btn{Remove transparent objects}, \btn{Break
        Subpaths}, and \btn{Reduce polyline details} in the optimisation
        settings. A drawing with 40\,000 nodes will never feel quick.
  \item Display: turn off \emph{Instant Snap Display}, animation and image
        caching toggles in \menu{View>Scene Appearance}, and \emph{Simplify
        Effects} (or hold \key{CapsLock}).
  \item Raster jobs: lower the DPI. Time scales linearly with it.
  \item Still slow: run once with \cmd{-f profile.txt} and send the result with a
        bug report --- it is a Python profile of the session, and it points at the
        actual bottleneck.
\end{itemize}

\subsection*{Something crashed}

MeerK40t installs a crash handler. On a crash it:

\begin{enumerate}[label=\arabic*.,leftmargin=1.4em]
  \item writes the traceback --- including a summary of the local variables at
        each frame --- to \file{MeerK40t-\emph{timestamp}.txt} in the working
        directory (falling back to the configuration directory if that is not
        writable), and to the terminal's standard error;
  \item shows a dialog titled \emph{Crash Detected! Send Log?} offering to send
        those details to the developers. The dialog is explicit that only the
        crash details are sent --- no project data, no personal information ---
        and \btn{Cancel} aborts the program rather than continuing.
\end{enumerate}

\noindent The crash log is the most useful thing you can attach to a report, and
it survives the crash: find it, read the first traceback, and note the version
from the title bar or \cmd{meerk40t --version} (the suffix, \cmd{git},
\cmd{src} or \cmd{pkg}, tells the developers how your copy was built). More on
reporting in Appendix~\ref{app:glossary}.

\subsection*{Permissions, sudo and USB}

If MeerK40t only sees the laser when run as root or with \cmd{sudo}, install the
udev rules from Chapter~\ref{ch:connect} and add your user to the \cmd{dialout}
group. Running the program with elevated privileges is not recommended, and the
program warns about it --- the warning can be silenced, but the better fix is the
rules.

\subsection*{Two copies of the program}

With \btn{Single Instance} on, launching a second copy brings the first one
forward. If you really do want two, or if the check is getting in your way on a
shared machine, turn it off in the \tabname{General} tab.

\begin{tryit}{Make your setup survivable}
Twenty minutes now, an afternoon saved later.
\begin{enumerate}[label=\arabic*.,leftmargin=1.4em]
  \item In the console, run \cmd{setting\_export /home/me/meerk40t-backup} and
        confirm the timestamped \file{.cfg} file is there. Do it again after any
        significant device or material work.
  \item Note your configuration file's location from the table above, and confirm
        the \file{.bak} files beside it.
  \item Set \emph{Levels of Undo-States} to something larger than the default if
        you do exploratory editing, and consider whether \emph{Autosave} interval
        of five minutes suits your machine.
  \item Turn on \emph{Enable Laser Arm} if anyone else ever uses this machine;
        leave it off if you are the only operator and find it annoying.
  \item Deliberately break the layout --- hide panes, float windows, move
        everything --- then fix it with \menu{Panes>Reset Panes} and
        \menu{Tools>Reset Windows}.
  \item Open a project, make a change, and confirm \btn{Save Window Positions}
        and \emph{Remember main window position} bring back your workspace on the
        next start.
\end{enumerate}
\end{tryit}
